% Folder 134-2, batch 10 — pages 181–184 (position in the folder; see below), the last
% four pages of the folder.
% Pass: Opus 5.5 (claude-opus-5-5), 2026-10-15 — first pass, unchecked against the pages by a human.
%
% Special rules for folder 134-2 (issues #44 and #57, the latter as decided by the
% user on 2026-10-15) applied: his own unprinted letters are transcribed in full.
%
% Numbering: \page{} numbers the position in the folder, as in batches 6 and 7 (and
% as the tiles and the facsimile do). From p. 115 on the archivists' pencil runs one
% ahead of the position: here p. 181 is pencilled 182, p. 184 is pencilled 185.
%
% Contents of the batch:
%  - pp. 181-184 (his pp. 37-40): end of the HANDWRITTEN FRENCH ORIGINAL of the
%    letter to Lawrence Breen, « Villecun 17/19 juillet 1975 », begun on p. 145 in
%    batch 8. The text starts mid-sentence (« les "complétés profinis" … de tels
%    espaces ») and ends on p. 184 with « Bien cordialement à toi / Alexandre
%    Grothendieck ». Transcribed in full, with the apparatus.
%
% Decisions: the English typed version (batch 4, pp. 75-77) and R. Brown's
% translation (batches 7-8) were used only to follow the argument, never to fill an
% \ill{}; where the English agrees with a doubtful reading, the note says so and the
% reading stays \uncertain{}. The struck parenthesis on Quillen at the top of the
% English p. 77 (« N.B. A first approach … ») has no counterpart in this French
% original. His script P for the stratification is set $\mathcal{P}$.
%
% Letters in this batch: Grothendieck to L. Breen, 17/19.7.1975 (end), French,
% handwritten, signed. No other document.
\documentclass[11pt,a4paper]{article}
\input{../preamble/grothendieck}

\folder{134-2}
\batch{10}
\pages{181}{184}
\foldertitle{[Chapitre I : Take off (pages 1 à 65) et table des matières provisoire] : tapuscrits annotés (19/02-22/02/1983), note manuscrite (s.d.), copies de lettre (1975, s.d.).}
\dating{1975-[1983]}
\watermark{Édition de démonstration}

\begin{document}

\section{Lettre à Lawrence Breen, Villecun, 17/19 juillet 1975 (original manuscrit, fin)}

\page{181}
\note{numérotation : à partir de la p.~115, le crayon des archivistes a un de plus que la position dans le dossier ; cette page est crayonnée 182. Pagination de la lettre : « 37 ». Suite de l'original français manuscrit de la lettre à L. Breen du 17/19 juillet 1975, commencé p.~145 ; la phrase en cours se poursuit ici (dans la version anglaise dactylographiée, lot~4, p.~75 : « their profinite completions … are those of such polyhedra »).}
les « complétés profinis » \add{(au sens d'Artin--Mazur),} premiers à $p$ (la car.\ de $k$) de tels espaces. On voit clairement comment \uncertain{aborder} un tel programme en car.~$0$, mais on prévoit de l'amusement supplémentaire, voire du mystère, en car.\ $p > 0$, pour les variétés qui, même birationnellement, refuseraient de se laisser remonter en car.~$0$ !
\note{« (au sens d'Artin--Mazur), » est écrit au-dessus de la première ligne, un trait le rattachant après « profinis ». « aborder » : le début du mot est mal formé ; l'anglais dit « how to begin on ».}

De ces cogitations essentiellement géométriques, je ne saurais à l'heure actuelle tirer un programme précis de développement de structures algébriques adéquates pour les exprimer. Je me bornerai à quelques remarques en marge.
\note{l'alinéa est précédé d'un trait horizontal qui sépare les deux parties.}

\struck{\ill{}} Depuis longtemps j'ai été intrigué par l'idée d'une « linéarisation » d'un type d'homotopie (ordinaire), i.e.\ la question du type : Si $X$ est un type d'homotopie, combien d'information cohomologique du type : cohomologie de $X$ à coefficients variables $M$ (constants ou constants tordus), structure multiplicative $H^i(X, M) \times H^j(X, N) \to H^{i+j}(X, M \otimes N)$, plus éventuellement d'autres opérations cohomologiques -- est-il nécessaire d'avoir pour reconstituer le type d'homotopie entièrement ? (disons,
\note{devant « Depuis », un signe court noirci, suivi d'une parenthèse fermante : peut-être un appel biffé, illisible. Le point d'interrogation après « entièrement » est surchargé sur un autre signe.}

\page{182}
\note{crayonnée 183. Pagination de la lettre : « 38 ».}
dans cette première approche pré-catégorie dérivée, par la connaissance du groupe fondamental $\pi_1$, donc de la catégorie des coeff.\ constants tordus $= \pi_1$-modules, les foncteurs $H^i(-, M)$ sur ceux-ci avec \uncertain{la} structure de foncteurs \struck{\ill{}, de \uncertain{foncteurs}} \add{cohomologiques} relatif aux suites exactes, la structure de cup-produit etc. -- reliés par certaines propriétés formelles.
\note{« cohomologiques » est écrit au-dessus des mots biffés. « \uncertain{la} structure » : l'article se lit plutôt « le », et le début de « structure » est surchargé ; l'anglais dit « the structure of cohomological functors ». « relatif » : sic.}

Une fois qu'on dispose du langage des cat.\ dérivées, la \add{sous-catégorie de la catégorie dérivée des} complexes abéliens, formée \add{des complexes de} faisceaux de coh.\ loc.\ constants sur $X$, avec sa structure triangulée et sa structure multiplicative $\overset{\mathbb{L}}{\otimes}$ (et éventuellement $\mathbb{R}\mathrm{Hom}$ …) \uncertain{donnant} un candidat plus satisfaisant pour espérer récupérer le type d'homotopie. Je n'ai guère d'idée si oui ou non ils permettent de le récupérer \add{en effet}, mais par contre n'ai guère de doute qu'en allant au bout de la « linéarisation », c'est à dire en passant dans le \emph{non abélien} et travaillant avec la $(n+1)$-catégorie (sans structure supplémentaire dessus !) des $n$-champs loc.\ constants \struck{\ill{}} pour tout $n$
\note{les deux insertions sont écrites au-dessus des lignes et soulignées d'un trait qui les rattache. « donnant » : lecture du tracé, la syntaxe attendrait « donne ». « par contre n'ai guère » : le sujet « je » manque sur la page. « loc.\ constants » est d'une encre plus appuyée, à la dernière ligne ; il est suivi de quelques lettres biffées et d'une croix, illisibles. L'anglais (lot~4, p.~76) porte ici « of constructible sheaves on $X$ », que cet original ne porte pas.}

\page{183}
\note{crayonnée 184. Pagination de la lettre : « 39 ».}
on arrive à reconstituer le type d'homotopie via son $\infty$-groupoïde fondamental, comme expliqué dans ma lettre précédente, et rappelé dans celle-ci. (Cela signifie en particulier que toutes les opérations cohomologiques possibles et imaginables sont déjà incluses dans la donnée d'un tel système de $n$-catégories ….) \struck{Mais}
\note{« et » (avant « rappelé ») est surchargé.}

De même, les types d'homotopie plus \uncertain{cumulés}, liés à certains diagrammes finis, qu'on peut associer à certains types de stratifications \add{$\mathcal{P}$} d'espaces topologiques « modérés » \add{$X$} disons, devraient correspondre de façon aussi parfaite aux \add{$(n+1)$-}catégories de $n$-champs sur $X$ qui sont loc.\ constants sur chacune des strates de $\mathcal{P}$ (disons : qui sont subordonnés à $\mathcal{P}$). Si plus haut la \ill{} de la cat.\ dérivée loc.\ constante marchait bel et bien, on s'attendrait ici à récupérer le type d'homotopie mixte à partir de la \add{s-}catégorie de la catégorie dérivée des faisceaux abéliens sur $X$, formée des complexes qui ont \uncertain{coh.}\ loc.\ constant sur chacune des strates -- avec en plus les opérations $\overset{\mathbb{L}}{\otimes}$, $\mathbb{R}\mathrm{Hom}$ plus, le
\note{« cumulés » : le mot commence par un trait qui ressemble à une parenthèse ouvrante ; l'anglais dit « more elaborate ». « certains types » : « certains » est surchargé. « $\mathcal{P}$ » est écrit au-dessus de « stratifications », « $X$ » ajouté dans l'interligne après « modérés », « $(n+1)$- » au-dessus de « catégories », où « aux » est surchargé. Le mot illisible, de quatre ou cinq lettres (« to?is »), désigne ce qui, plus haut, « marchait » ; l'anglais dit « description of homotopy types by the « locally constant derived category » ». « s- » au-dessus de « catégorie » abrège « sous- ». « \uncertain{coh.} » : précédé d'un signe en deux points. « $\mathbb{R}\mathrm{Hom}$ » est souligné.}

\page{184}
\note{crayonnée 185. Pagination de la lettre : « 40 ».}
cas échéant, les $4$ opérations $\mathbb{R}g_!$, $\mathbb{R}g_*$, $\mathbb{L}g^*$, $g^!$ pour des \uncertain{induits} $g \colon Z' \to Z''$ de réunions localement fermées de strates …. (L'ennui ici est qu'on ne sait même pas encore ce que c'est qu'une catégorie triangulée, pas plus qu'on ne sait ce qu'est sa version non commutative, sans doute plus simple et plus fondamentale : \emph{une} « catégorie homotopique », avec opérations « fibre » et « cofibre » ….)
\note{« \uncertain{induits} » : l'anglais dit « for the induced $g$ », ce qui s'accorde avec la lecture, sans la trancher. Le $L$ de « $\mathbb{L}g^*$ » est tracé simple.}

Il est vraiment temps que je termine cette lettre-fleuve, qui devient de plus en plus vaseuse. Juste une question : Quelle est cette mirifique formule de Bloch--Quillen à laquelle tu fais allusion, dont je n'ai jamais entendu parler, et qui fait venir l'eau à la bouche ?
\note{l'alinéa est précédé d'un trait horizontal. La fin de « Juste » est surchargée.}

Bien cordialement à toi

Alexandre

Grothendieck
\note{fin de la lettre à L. Breen du 17/19 juillet 1975, et du dossier.}

\end{document}
